\documentclass{article}
\usepackage[utf8]{inputenc}
\usepackage{amsmath}
\usepackage{geometry}
\geometry{margin=2cm}
\begin{document}

\section*{Questão 115 -- ENEM 2024 -- Ciências da Natureza}

\subsection*{Tema: Estequiometria e concentrações de soluções}

\subsubsection*{Análise e Interpretação da Questão}
A questão solicita o cálculo da massa de cloreto de sódio (NaCl) contida em uma colher pequena da receita de soro caseiro,
a partir da concentração de \i{íon sódio} e das quantidades de ingredientes dadas.

\paragraph{Estratégia para resolução} Deve-se calcular inicialmente a massa total de sódio na solução,
utilizando a concentração e o volume total. Em seguida, utilizando as massas molares de Na e NaCl,
calcula-se a massa total do sal. Por fim, divide-se essa massa total pelo número de colheres pequenas da receita
(2 colheres), para encontrar a massa por colher.

\subsection*{Conteúdos Envolvidos e Mini Revisão}

\begin{tabular}{|p{5cm}|p{8cm}|}
\hline
\textbf{Conceito} & \textbf{Aplicacão} \\
\hline
Concentração e volume & Multiplicação da concentração pela quantidade total para obter massa do \i{ion}. \\
\hline
Estequiometria & Uso da proporção entre massas molares para relacionar massa do \i{íon} \`a do composto molecular. \\
\hline
Divisão proporcional & Divisão da massa total pelo número de colheres para obter massa por unidade. \\
\hline
\end{tabular}

\paragraph{Teorias Relevantes}
\begin{itemize}
  \item Concentração (mg/mL) indica a quantidade de soluto por mililitro.
  \item Estequiometria relaciona as quantidades de elementos e compostos químicos baseados em massas molares.
  \item Divisão proporcional \`e necessária para calcular a massa contida em unidades parciais de receitu\a.
\end{itemize}

\subsection*{Resolução Detalhada e Pulo do Gato}

\begin{enumerate}
  \item Converte-se o volume total: $2\,\mathrm{L} = 2000\,\mathrm{mL}$.
  \item Calcula-se a massa de \i{sódio} total: $1,4 \frac{\mathrm{mg}}{\mathrm{mL}} \times 2000\,\mathrm{mL} = 2800\,\mathrm{mg} = 2,8\,\mathrm{g}$.
  \item Determina-se a razão entre massa molar de NaCl e Na: $\frac{58,5}{23} \approx 2,54$.
  \item Multiplica-se para achar massa de NaCl: $2,8\,\mathrm{g} \times 2,54 = 7,1\,\mathrm{g}$.
  \item Divide-se pela quantidade de colheres pequenas (2): $7,1\,\mathrm{g} \div 2 = 3,55\,\mathrm{g}$.
  \item Alternativa próxima: letra D, 3,6 g.
\end{enumerate}

\paragraph{Pulo do gato} Fazer o cálculo primeiro da massa do sódio total e usar a proporção das massas molares para considerar o sal inteiro, dividindo corretamente pelo número de colheres.

\subsection*{Comentários sobre as Alternativas}

\begin{itemize}
  \item \textbf{A}: Valor considerado apenas com massa de sódio, sem incluir cloro - incorreto.
  \item \textbf{B}: Considera massa de sódio, mas não converte para NaCl ou erra na divisão.
  \item \textbf{C}: Massa total correta, mas divide erradamente pelo número de colheres.
  \item \textbf{D}: Correta, aplica corretamente todos os cálculos e divisões.
  \item \textbf{E}: Confunde massa total para duas colheres como massa para uma só, superestimando o valor.
\end{itemize}

\subsection*{Código da Questão}
CN_EST_SOL_001

\bigskip
\noindent\textbf{\large Em resumo, a resolução exige calculação cuidadosa da massa total do \i{ion} e a utilização correta das proporções molares, seguido da divisão pela quantidade correta de unidades da medida da receita, garantindo um resultado coerente e preciso.}

\end{document}